\documentclass[12pt,a4paper]{article}
\usepackage[utf8]{vietnam}
\usepackage{amsmath,amssymb}
\usepackage{tabularx}
\usepackage{tikz}
\usepackage{xcolor}
% Giảm lề tối đa tận dụng diện tích giấy (chuẩn máy in laser/photo)
\usepackage[top=0.6cm, bottom=0.6cm, left=0.6cm, right=0.6cm]{geometry}
\usepackage[dethi]{ex_test}

\setlength{\parskip}{0pt}

% Dòng kẻ chấm mờ đều đặn chuẩn viết tay (khoảng cách 7.2mm rộng rãi)
\newcommand{\dongke}[1]{%
  \begingroup%
  \setlength{\parskip}{0pt}%
  \baselineskip=20.5pt%
  \vspace{1.5mm}%
  \foreach \i in {1,...,#1}{%
    \noindent{\color{gray!60}\dotfill}\par%
  }%
  \vspace{1mm}%
  \endgroup%
}

% Định nghĩa ô bài tập tự luận trong bảng
\newcommand{\oBaiTap}[3][7]{%
  \vspace{1.5mm}%
  \noindent\textbf{#2} #3\par%
  \dongke{#1}%
}

% Câu trắc nghiệm 1 cột: 4 đáp án dàn hàng ngang trên 1 dòng
\newcommand{\cauTN}[6]{% \cauTN{số câu}{đề}{A}{B}{C}{D}
  \noindent\textbf{Câu #1.} #2\par\vspace{0.4mm}%
  \noindent%
  \makebox[0.25\linewidth][l]{\textbf{A.} #3}%
  \makebox[0.25\linewidth][l]{\textbf{B.} #4}%
  \makebox[0.25\linewidth][l]{\textbf{C.} #5}%
  \makebox[0.25\linewidth][l]{\textbf{D.} #6}\par\vspace{0.8mm}%
}

\begin{document}
\pagestyle{empty}

% #############################################################################
% ========================= TOÁN 6 - TUẦN 1 - BUỔI 1: MÃ ĐỀ 101 =============
% #############################################################################

% -------------------- TRANG 1 (MÃ ĐỀ 101): TRẮC NGHIỆM --------------------
\noindent
\begin{tabularx}{\linewidth}{|X|c|}
\hline
\textbf{BÀI 1: TẬP HỢP (TUẦN 1 - BUỔI 1)} & \textbf{Mã đề: 101} \\
\hline
\end{tabularx}

\vspace{1.2mm}

\noindent
\begin{tabularx}{\linewidth}{|X|}
\hline
\rule[-3mm]{0pt}{9.5mm}\textbf{Họ và tên:} \dotfill \quad \textbf{Lớp:} \makebox[1.8cm]{\dotfill} \\
\hline
\end{tabularx}

\vspace{1.2mm}

\noindent
\begin{tabularx}{\linewidth}{|p{3cm}|X|}
\hline
\centering\textbf{ĐIỂM SỐ} & \centering\textbf{LỜI NHẬN XÉT} \tabularnewline
\hline
\rule{0pt}{2.2cm} & \tabularnewline
\hline
\end{tabularx}

\vspace{1.5mm}
\noindent\textbf{BẢNG ĐIỀN ĐÁP ÁN TRẮC NGHIỆM (Mỗi câu đúng 0,25 điểm):}
\vspace{1.2mm}

\noindent
\renewcommand{\arraystretch}{1.5}
\begin{tabularx}{\linewidth}{|c|*{12}{>{\centering\arraybackslash}X|}}
\hline
\textbf{Câu} & \textbf{1} & \textbf{2} & \textbf{3} & \textbf{4} & \textbf{5} & \textbf{6} & \textbf{7} & \textbf{8} & \textbf{9} & \textbf{10} & \textbf{11} & \textbf{12} \\
\hline
\textbf{Đ/A} & & & & & & & & & & & & \\
\hline
\end{tabularx}
\renewcommand{\arraystretch}{1.0}

\vspace{3mm}

\noindent\textbf{\large PHẦN I. TRẮC NGHIỆM KHÁCH QUAN (3,0 điểm)}
\vspace{1.5mm}

\cauTN{1}{Cho tập hợp $A = \{2; 3; 5; 7\}$. Khẳng định nào sau đây là \textbf{đúng}?}{$2 \notin A$}{$7 \notin A$}{$5 \in A$}{$3 \subset A$}

\cauTN{2}{Kí hiệu nào sau đây dùng để chỉ phần tử $x$ không thuộc tập hợp $M$?}{$x \in M$}{$x \notin M$}{$x \subset M$}{$M \in x$}

\cauTN{3}{Viết tập hợp $B$ các chữ cái trong từ ``TOAN'' bằng cách liệt kê phần tử:}{$B = \{\text{T}; \text{O}; \text{A}; \text{N}\}$}{$B = \{\text{T}; \text{O}; \text{N}\}$}{$B = \{\text{T}; \text{A}; \text{N}\}$}{$B = \{\text{O}; \text{A}; \text{N}\}$}

\cauTN{4}{Cho tập hợp $C = \{x \in \mathbb{N} \mid 3 < x \le 7\}$. Dạng liệt kê của tập hợp $C$ là:}{$C = \{3; 4; 5; 6; 7\}$}{$C = \{4; 5; 6\}$}{$C = \{4; 5; 6; 7\}$}{$C = \{3; 4; 5; 6\}$}

\cauTN{5}{Số phần tử của tập hợp $D = \{10; 11; 12; \dots; 25\}$ là:}{$15$}{$16$}{$14$}{$25$}

\cauTN{6}{Tập hợp các số tự nhiên chẵn nhỏ hơn 10 là:}{$\{0; 2; 4; 6; 8\}$}{$\{2; 4; 6; 8\}$}{$\{0; 2; 4; 6; 8; 10\}$}{$\{2; 4; 6; 8; 10\}$}

\cauTN{7}{Tập hợp các số tự nhiên khác 0 được kí hiệu là:}{$\mathbb{N}$}{$\mathbb{N}^*$}{$\mathbb{Z}$}{$\mathbb{Z}^*$}

\cauTN{8}{Số phần tử của tập hợp $E = \{5; 10; 15; 20; \dots; 50\}$ là:}{$10$}{$9$}{$11$}{$8$}

\cauTN{9}{Trong các khẳng định sau, khẳng định nào \textbf{sai}?}{$0 \in \mathbb{N}$}{$0 \in \mathbb{N}^*$}{$1 \in \mathbb{N}^*$}{$5 \in \mathbb{N}$}

\cauTN{10}{Tập hợp các chữ cái trong từ ``HOC SINH'' gồm bao nhiêu phần tử?}{$7$}{$6$}{$5$}{$8$}

\cauTN{11}{Cho tập hợp $M = \{x \in \mathbb{N} \mid x + 5 = 12\}$. Số phần tử của tập hợp $M$ là:}{$0$}{$1$}{$2$}{$7$}

\cauTN{12}{Tập hợp các số tự nhiên $x$ sao cho $x \cdot 0 = 0$ là:}{$\emptyset$}{$\{0\}$}{$\mathbb{N}$}{$\mathbb{N}^*$}

\newpage
% -------------------- TRANG 2 (MÃ ĐỀ 101): TỰ LUẬN --------------------
\noindent
\begin{tabularx}{\linewidth}{|X|c|}
\hline
\textbf{PHẦN II. TỰ LUẬN (7,0 ĐIỂM)} & \textbf{Mã đề: 101} \\
\hline
\end{tabularx}

\vspace{1.2mm}

\noindent
\begin{tabularx}{\linewidth}{|X|}
\hline
\rule[-3mm]{0pt}{9.5mm}\textbf{Họ và tên:} \dotfill \quad \textbf{Lớp:} \makebox[1.8cm]{\dotfill} \\
\hline
\end{tabularx}

\vspace{2.5mm}

\noindent\textbf{\large Câu 1 (2,5 điểm). Viết và mô tả tập hợp:}
\vspace{1.2mm}

\noindent
\begin{tabularx}{\linewidth}{|X|X|}
\hline
\oBaiTap[7]{a)}{Viết tập hợp $A$ các số tự nhiên lớn hơn $7$ và không vượt quá $13$ bằng hai cách (liệt kê các phần tử và chỉ ra tính chất đặc trưng).} & 
\oBaiTap[7]{b)}{Viết tập hợp $B$ các chữ cái trong cụm từ ``NHA TRANG''. Tập hợp $B$ có tất cả bao nhiêu phần tử?} \\
\hline
\end{tabularx}

\vspace{3.5mm}

\noindent\textbf{\large Câu 2 (2,5 điểm). Số phần tử và Tìm phần tử của tập hợp:}
\vspace{1.2mm}

\noindent
\begin{tabularx}{\linewidth}{|X|X|}
\hline
\oBaiTap[7]{a)}{Tính số phần tử của tập hợp các số tự nhiên: $C = \{12; 15; 18; \dots; 96\}$.} & 
\oBaiTap[7]{b)}{Cho tập hợp $D = \{x \in \mathbb{N} \mid (x - 3) \cdot 4 = 24\}$. Hãy tìm phần tử $x$ và viết lại tập hợp $D$.} \\
\hline
\end{tabularx}

\vspace{3.5mm}

\noindent\textbf{\large Câu 3 (2,0 điểm). Bài toán thực tế:}
\vspace{1.2mm}

\noindent
\begin{tabularx}{\linewidth}{|X|}
\hline
\vspace{0.8mm}
Một lớp học có $42$ học sinh. Trong đó có $26$ bạn thích môn Bóng đá, $18$ bạn thích môn Cầu lông, và có $8$ bạn thích cả hai môn thể thao đó. Biết rằng trong lớp bạn nào cũng thích ít nhất một trong hai môn. Hỏi lớp học đó có bao nhiêu học sinh chỉ thích môn Bóng đá và bao nhiêu học sinh chỉ thích môn Cầu lông?
\par\vspace{1.5mm}
\noindent\textit{Bài giải:}
\dongke{8}
\\ \hline
\end{tabularx}

\newpage
% #############################################################################
% ========================= TOÁN 6 - TUẦN 1 - BUỔI 1: MÃ ĐỀ 102 =============
% #############################################################################

% -------------------- TRANG 3 (MÃ ĐỀ 102): TRẮC NGHIỆM --------------------
\noindent
\begin{tabularx}{\linewidth}{|X|c|}
\hline
\textbf{BÀI 1: TẬP HỢP (TUẦN 1 - BUỔI 1)} & \textbf{Mã đề: 102} \\
\hline
\end{tabularx}

\vspace{1.2mm}

\noindent
\begin{tabularx}{\linewidth}{|X|}
\hline
\rule[-3mm]{0pt}{9.5mm}\textbf{Họ và tên:} \dotfill \quad \textbf{Lớp:} \makebox[1.8cm]{\dotfill} \\
\hline
\end{tabularx}

\vspace{1.2mm}

\noindent
\begin{tabularx}{\linewidth}{|p{3cm}|X|}
\hline
\centering\textbf{ĐIỂM SỐ} & \centering\textbf{LỜI NHẬN XÉT} \tabularnewline
\hline
\rule{0pt}{2.2cm} & \tabularnewline
\hline
\end{tabularx}

\vspace{1.5mm}
\noindent\textbf{BẢNG ĐIỀN ĐÁP ÁN TRẮC NGHIỆM (Mỗi câu đúng 0,25 điểm):}
\vspace{1.2mm}

\noindent
\renewcommand{\arraystretch}{1.5}
\begin{tabularx}{\linewidth}{|c|*{12}{>{\centering\arraybackslash}X|}}
\hline
\textbf{Câu} & \textbf{1} & \textbf{2} & \textbf{3} & \textbf{4} & \textbf{5} & \textbf{6} & \textbf{7} & \textbf{8} & \textbf{9} & \textbf{10} & \textbf{11} & \textbf{12} \\
\hline
\textbf{Đ/A} & & & & & & & & & & & & \\
\hline
\end{tabularx}
\renewcommand{\arraystretch}{1.0}

\vspace{3mm}

\noindent\textbf{\large PHẦN I. TRẮC NGHIỆM KHÁCH QUAN (3,0 điểm)}
\vspace{1.5mm}

\cauTN{1}{Cho tập hợp $M = \{1; 4; 6; 9\}$. Khẳng định nào sau đây là \textbf{đúng}?}{$1 \notin M$}{$6 \in M$}{$4 \subset M$}{$9 \notin M$}

\cauTN{2}{Kí hiệu nào sau đây dùng để chỉ phần tử $a$ thuộc tập hợp $S$?}{$a \notin S$}{$a \subset S$}{$a \in S$}{$S \in a$}

\cauTN{3}{Viết tập hợp $A$ các chữ cái trong từ ``SACH'' bằng cách liệt kê phần tử:}{$A = \{\text{S}; \text{A}; \text{C}; \text{H}\}$}{$A = \{\text{S}; \text{A}; \text{C}\}$}{$A = \{\text{A}; \text{C}; \text{H}\}$}{$A = \{\text{S}; \text{C}; \text{H}\}$}

\cauTN{4}{Cho tập hợp $B = \{x \in \mathbb{N} \mid 5 \le x < 9\}$. Dạng liệt kê của tập hợp $B$ là:}{$B = \{5; 6; 7; 8; 9\}$}{$B = \{6; 7; 8\}$}{$B = \{5; 6; 7; 8\}$}{$B = \{6; 7; 8; 9\}$}

\cauTN{5}{Số phần tử của tập hợp $C = \{20; 21; 22; \dots; 35\}$ là:}{$15$}{$16$}{$14$}{$35$}

\cauTN{6}{Tập hợp các số tự nhiên lẻ nhỏ hơn 10 là:}{$\{1; 3; 5; 7; 9\}$}{$\{3; 5; 7; 9\}$}{$\{1; 3; 5; 7\}$}{$\{0; 1; 3; 5; 7; 9\}$}

\cauTN{7}{Tập hợp các số tự nhiên kí hiệu là $\mathbb{N}$. Phần tử nào sau đây thuộc $\mathbb{N}^*$?}{$0$}{$1$}{$-1$}{$\dfrac{1}{2}$}

\cauTN{8}{Số phần tử của tập hợp $K = \{4; 8; 12; 16; \dots; 48\}$ là:}{$11$}{$12$}{$13$}{$10$}

\cauTN{9}{Trong các khẳng định sau, khẳng định nào \textbf{đúng}?}{$0 \in \mathbb{N}^*$}{$0 \notin \mathbb{N}$}{$0 \in \mathbb{N}$}{$-2 \in \mathbb{N}$}

\cauTN{10}{Tập hợp các chữ cái trong từ ``HOC TOAN'' gồm bao nhiêu phần tử?}{$7$}{$6$}{$8$}{$5$}

\cauTN{11}{Cho tập hợp $P = \{x \in \mathbb{N} \mid x - 4 = 11\}$. Số phần tử của tập hợp $P$ là:}{$1$}{$0$}{$2$}{$15$}

\cauTN{12}{Tập hợp các số tự nhiên $x$ sao cho $x \cdot 0 = 5$ là:}{$\mathbb{N}$}{$\{0\}$}{$\emptyset$}{$\{5\}$}

\newpage
% -------------------- TRANG 4 (MÃ ĐỀ 102): TỰ LUẬN --------------------
\noindent
\begin{tabularx}{\linewidth}{|X|c|}
\hline
\textbf{PHẦN II. TỰ LUẬN (7,0 ĐIỂM)} & \textbf{Mã đề: 102} \\
\hline
\end{tabularx}

\vspace{1.2mm}

\noindent
\begin{tabularx}{\linewidth}{|X|}
\hline
\rule[-3mm]{0pt}{9.5mm}\textbf{Họ và tên:} \dotfill \quad \textbf{Lớp:} \makebox[1.8cm]{\dotfill} \\
\hline
\end{tabularx}

\vspace{2.5mm}

\noindent\textbf{\large Câu 1 (2,5 điểm). Viết và mô tả tập hợp:}
\vspace{1.2mm}

\noindent
\begin{tabularx}{\linewidth}{|X|X|}
\hline
\oBaiTap[7]{a)}{Viết tập hợp $M$ các số tự nhiên lớn hơn $5$ và không vượt quá $11$ bằng hai cách (liệt kê các phần tử và chỉ ra tính chất đặc trưng).} & 
\oBaiTap[7]{b)}{Viết tập hợp $N$ các chữ cái trong cụm từ ``DA LAT''. Tập hợp $N$ có tất cả bao nhiêu phần tử?} \\
\hline
\end{tabularx}

\vspace{3.5mm}

\noindent\textbf{\large Câu 2 (2,5 điểm). Số phần tử và Tìm phần tử của tập hợp:}
\vspace{1.2mm}

\noindent
\begin{tabularx}{\linewidth}{|X|X|}
\hline
\oBaiTap[7]{a)}{Tính số phần tử của tập hợp các số tự nhiên: $P = \{14; 18; 22; \dots; 98\}$.} & 
\oBaiTap[7]{b)}{Cho tập hợp $Q = \{x \in \mathbb{N} \mid (x + 2) \cdot 5 = 35\}$. Hãy tìm phần tử $x$ và viết lại tập hợp $Q$.} \\
\hline
\end{tabularx}

\vspace{3.5mm}

\noindent\textbf{\large Câu 3 (2,0 điểm). Bài toán thực tế:}
\vspace{1.2mm}

\noindent
\begin{tabularx}{\linewidth}{|X|}
\hline
\vspace{0.8mm}
Một lớp học có $45$ học sinh. Trong đó có $28$ bạn thích môn Bơi lội, $20$ bạn thích môn Cờ vua, và có $9$ bạn thích cả hai môn đó. Biết rằng trong lớp bạn nào cũng thích ít nhất một trong hai môn. Hỏi lớp học đó có bao nhiêu học sinh chỉ thích môn Bơi lội và bao nhiêu học sinh chỉ thích môn Cờ vua?
\par\vspace{1.5mm}
\noindent\textit{Bài giải:}
\dongke{8}
\\ \hline
\end{tabularx}

\end{document}
